% ----------------------------------------------------------------------
%  2.1 (вводная часть) и 2.1.1. Вакансии и межузельные атомы
% ----------------------------------------------------------------------
\begin{orientir}
Точечные дефекты~— самые простые нарушения решётки размером в один-два
атома. Несмотря на малый размер, они управляют диффузией, во многом
определяют фазовые превращения и поведение материала под облучением, а
взаимодействуя с дислокациями, влияют на прочность. В разделе~\ref{sec:point}
мы выясним: какие бывают точечные дефекты (2.1.1); сколько энергии нужно,
чтобы их создать, сколько их в равновесии и как они движутся (2.1.2);
как они объединяются в комплексы (2.1.3); где рождаются и исчезают (2.1.4);
как получить их много (2.1.5) и как их измерить (2.1.6).
\end{orientir}

\subsection{Вакансии и межузельные атомы}\label{ssec:vac-int}

В чистом металле возможны всего два простейших точечных дефекта
(рис.~\ref{fig:vac-int-2d}).

\begin{opredelenie}
\term{Вакансия}~— узел кристаллической решётки, в котором нет атома.
\term{Межузельный атом}~— «лишний» атом основного вещества, находящийся
не в узле, а в промежутке между узлами (в~\term{междоузлии}).
\end{opredelenie}

\begin{masshtab}[точечный дефект]{-9.7/-9.2/{вакансия и зона искажений}, -8.55/-8.15/{между вакансиями у $T_\text{пл}$}, -2.6/-2.2/{одна вакансия при 300\,К}}
Вакансия~— пустой узел (0,25~нм); заметные смещения соседей затухают
уже через 1--2 межатомных расстояния. У температуры плавления в меди
одна вакансия приходится на $\sim10^4$ узлов, т.\,е. вакансии отстоят друг
от друга в среднем на $\sqrt[3]{10^4}\approx20$ межатомных расстояний
($\approx5$~нм). При комнатной температуре равновесная концентрация
$\sim10^{-21}$: одна вакансия на кубик металла размером в несколько
миллиметров (п.~\ref{ssec:formation}). Поэтому все вакансии, которые
удаётся наблюдать при комнатной температуре, неравновесные: они
закалены с высоких температур или созданы облучением и деформацией.
\end{masshtab}

В сплавах добавляются примесные (легирующие) атомы. Если примесный
атом занимает узел вместо атома основы, это \term{атом замещения}; если
он мал и располагается в междоузлии~— \term{атом внедрения} (так ведут
себя, например, C, N, H в железе).

\begin{figure}[ht]
\centering
\begin{tikzpicture}[scale=0.62]
  % смещения показаны с увеличением: к вакансии (2,2) и от межузельного атома (6.5,2.5)
  \foreach \x in {0,...,8}{\foreach \y in {0,...,4}{
    \pgfmathsetmacro{\dx}{0}\pgfmathsetmacro{\dy}{0}
    \pgfmathtruncatemacro{\code}{10*\x+\y}
    \ifnum\code=12 \pgfmathsetmacro{\dx}{0.2}\fi
    \ifnum\code=32 \pgfmathsetmacro{\dx}{-0.2}\fi
    \ifnum\code=21 \pgfmathsetmacro{\dy}{0.2}\fi
    \ifnum\code=23 \pgfmathsetmacro{\dy}{-0.2}\fi
    \ifnum\code=62 \pgfmathsetmacro{\dx}{-0.17}\pgfmathsetmacro{\dy}{-0.17}\fi
    \ifnum\code=72 \pgfmathsetmacro{\dx}{0.17}\pgfmathsetmacro{\dy}{-0.17}\fi
    \ifnum\code=63 \pgfmathsetmacro{\dx}{-0.17}\pgfmathsetmacro{\dy}{0.17}\fi
    \ifnum\code=73 \pgfmathsetmacro{\dx}{0.17}\pgfmathsetmacro{\dy}{0.17}\fi
    \ifnum\code=22 \node[vac] at (\x,\y) {};
    \else
      \pgfmathparse{abs(\dx)+abs(\dy)>0 ? 1 : 0}
      \ifnum\pgfmathresult=1
        \draw[atomLine!60, densely dotted, line width=0.4pt] (\x,\y) circle (0.34);
        \node[atom] at (\x+\dx,\y+\dy) {};
      \else \node[atom] at (\x,\y) {};\fi
    \fi}}
  \node[atomI] at (6.5,2.5) {};
  \node[lbl, below] at (2,-0.6) {вакансия};
  \node[lbl, below] at (6.5,-0.6) {межузельный атом};
\end{tikzpicture}
\caption{Вакансия и межузельный атом в двумерной модели кристалла.
Соседние атомы смещены (релаксация, смещения показаны с увеличением;
пунктир~— их положения в идеальной решётке): к вакансии~— внутрь,
от межузельного атома~— наружу}
\label{fig:vac-int-2d}
\end{figure}

\begin{glavnoe}
Положение вакансии однозначно~— это всегда узел решётки. Для
межузельного атома возможны \emph{разные конфигурации}, и энергия
кристалла зависит от того, какую из них атом выберет. Поэтому вопрос
«где именно сидит межузельный атом» нетривиален.
\end{glavnoe}

\subsubsection{Конфигурации межузельного атома в ГЦК-решётке}
Для ГЦК-решётки рассматривают три конфигурации (рис.~\ref{fig:int-fcc}).

\begin{enumerate}
\item \term{Объёмно-центрированная (октаэдрическая)}: атом сидит в центре
куба элементарной ячейки, в самой большой пустоте ГЦК-решётки.
\item \term{Гантельная (расщеплённая)}: межузельный атом подходит к узлу
в центре грани, выталкивает «хозяина» из узла, и оба атома располагаются
симметрично относительно этого узла вдоль направления $\nfam{100}$.
Получается пара~— «гантель». Расчёт для меди даёт расстояние между атомами
гантели около $0{,}6\,a$.
\item \term{Краудион} \eng{crowd~— толпа}: лишний атом «вдавлен» в
плотноупакованный ряд $\nfam{110}$. Атомы ряда сжаты, и это сжатие
плавно спадает в обе стороны на расстояние до $\sim 10\,a$.
\end{enumerate}

\begin{figure}[ht]
\centering
\begin{subfigure}[b]{0.3\textwidth}\centering
\begin{tikzpicture}[cubeproj, scale=2.25]
  \CubeEdges
  \FCCback \FCCmid
  \node[atomI] at (0.5,0.5,0.5) {};
  \FCCfront
\end{tikzpicture}
\caption{октаэдрическая}
\end{subfigure}\hfill
\begin{subfigure}[b]{0.3\textwidth}\centering
\begin{tikzpicture}[cubeproj, scale=2.25]
  \CubeEdges
  \FCCback \FCCmid
  \foreach \p in {(0,0,0),(1,0,0),(0,0,1),(1,0,1)} \node[atom=3.4mm] at \p {};
  \node[vac=3.4mm] at (0.5,0,0.5) {};
  \node[atomI] (g1) at (0.2,0,0.5) {};
  \node[atomI] (g2) at (0.8,0,0.5) {};
  \draw[dimline] (0.2,0,0.32) -- node[below, slbl]{$0{,}6a$} (0.8,0,0.32);
\end{tikzpicture}
\caption{гантель $\nfam{100}$}
\end{subfigure}\hfill
\begin{subfigure}[b]{0.36\textwidth}\centering
\begin{tikzpicture}[scale=0.42]
  % треугольная (плотноупакованная) сетка, средний ряд — краудион
  \foreach \r in {0,2}{\foreach \k in {-5,...,5}{
     \node[atom=4.0mm] at (\k, {\r*0.866}) {};}}
  \foreach \k in {-6,...,6}{
     \pgfmathsetmacro{\xk}{\k - 0.5*tanh(\k/2.2)}
     \pgfmathsetmacro{\shade}{100-55*exp(-(\k/3)^2)}
     \node[atom=4.0mm, fill=atomI!\shade!atomA] at (\xk, 0.866) {};}
  \node[slbl, below] at (0,-0.5) {ряд $\nfam{110}$};
\end{tikzpicture}
\caption{краудион}
\end{subfigure}
\caption{Конфигурации межузельного атома в ГЦК-решётке.
\textit{а}~— атом в центре куба;
\textit{б}~— расщеплённая (гантельная) конфигурация: узел в центре
грани (пунктир) занят парой атомов;
\textit{в}~— краудион: в среднем плотноупакованном ряду на один атом
больше, чем в соседних; сильнее всего сжата его середина (цвет)}
\label{fig:int-fcc}
\end{figure}

Какая из конфигураций устойчива, выяснили компьютерным моделированием.
Атом, помещённый в центр куба, «скатывается» к одному из ближайших узлов
и образует гантель~— значит, \vazhno{в ГЦК-металлах устойчива
гантельная конфигурация}. Однако разница энергий гантельной и
объёмно-центрированной конфигураций мала, порядка $0{,}2\eV$; это
важно для подвижности межузельного атома (с.~\pageref{par:int-mobility}).

В ОЦК-решётке (например, в $\alpha$-железе) устойчива гантель,
ориентированная вдоль $\nfam{110}$ (рис.~\ref{fig:int-bcc}).

\begin{figure}[ht]
\centering
\begin{tikzpicture}[cubeproj, scale=2.25]
  \CubeEdges
  \foreach \p in {(0,1,0),(1,1,0),(0,1,1),(1,1,1)} \node[atom=3.4mm] at \p {};
  \node[vac=3.4mm] at (0.5,0.5,0.5) {};
  \node[atomI] at (0.3,0.3,0.5) {};
  \node[atomI] at (0.7,0.7,0.5) {};
  \draw[fmRed, dashed, line width=0.5pt] (0.05,0.05,0.5) -- (0.95,0.95,0.5);
  \foreach \p in {(0,0,0),(1,0,0),(0,0,1),(1,0,1)} \node[atom=3.4mm] at \p {};
  \node[slbl, fmRed, right] at (0.95,0.95,0.5) {$\nfam{110}$};
\end{tikzpicture}
\caption{Гантельная конфигурация межузельного атома в ОЦК-решётке:
узел в центре куба (пунктир) занят парой атомов, ориентированной вдоль
$\nfam{110}$}
\label{fig:int-bcc}
\end{figure}

\begin{itogi}
\begin{itemize}
\item В чистом металле два собственных точечных дефекта: вакансия (пустой
узел) и межузельный атом (лишний атом в междоузлии).
\item В сплавах добавляются атомы замещения (в узлах) и внедрения
(в междоузлиях).
\item Межузельный атом может иметь несколько конфигураций. Устойчивая
в ГЦК~— гантель $\nfam{100}$, в ОЦК~— гантель $\nfam{110}$.
Краудион~— лишний атом в плотноупакованном ряду, вытянутый вдоль ряда.
\end{itemize}
\end{itogi}
